\documentclass[10pt,twocolumn]{article}
\usepackage[margin=1in]{geometry}
\usepackage{graphicx}
\usepackage{amsmath,amssymb,amsfonts}
\usepackage{algorithm}
\usepackage{algorithmic}
\usepackage{booktabs}
\usepackage{multirow}
\usepackage{hyperref}
\usepackage{natbib}
\usepackage{caption}
\usepackage{xcolor}
\usepackage{subcaption}

\title{Representation-Aware Multi-Objective Optimization\\for Multimodal Cancer Prognosis}
\author{Ankit Bhardwaj \\
University of Connecticut \\
\texttt{ankit.bhardwaj@uconn.edu}}
\date{}

\begin{document}

\maketitle

% ─────────────────────────────────────────────────────────────────────────────
\begin{abstract}
Multimodal cancer survival prediction integrates whole-slide
histopathology images with bulk RNA-seq transcriptomes, but joint
training under Cox partial likelihood faces a compounding challenge:
auxiliary objectives---such as per-modality survival heads and pairwise
ranking losses---can produce gradients that \textit{conflict} with the
primary objective's gradient direction in encoder parameter space,
degrading rather than improving the fused model. We introduce
\textbf{RA-SAMO} (Representation-Aware Survival-Aware Multi-Objective
optimization), a framework that (1)~decomposes training into a primary
Cox loss and multiple auxiliary objectives with adaptive EMA-based
weighting, and (2)~projects auxiliary encoder gradients that conflict
with the primary gradient onto its orthogonal complement, preserving
the primary objective's descent direction while retaining
non-conflicting auxiliary signal. We evaluate RA-SAMO on five TCGA
cohorts (BRCA, KIRC, LUAD, LUSC, STAD) under stratified 5-fold
cross-validation and demonstrate generalization to 7-class
dermoscopic lesion classification on HAM10000. Gradient-level
diagnostics---including Fisher information traces and inter-objective
cosine similarities---reveal that objective effectiveness is
cancer-dependent: ranking loss gradients align closely with Cox
($\cos > 0.93$) and thus add little, while auxiliary survival heads
provide orthogonal signal ($\cos \approx 0.14$) but require careful
weighting. RA-SAMO achieves competitive C-index across all five
cancers (mean 0.629) with the lowest cross-method variance, and on
HAM10000 produces the most stable macro-F1 ($0.695 \pm 0.011$) across
seeds. These results position RA-SAMO as a principled framework for
multi-objective multimodal optimization in data-scarce clinical settings,
with clear scope conditions: the framework provides the largest gains
when auxiliary gradients are both informative and conflicting.
\end{abstract}


% ─────────────────────────────────────────────────────────────────────────────
\section{Introduction}

Cancer prognosis is a critical clinical task that supports treatment
planning, patient stratification, and personalized therapeutic
decisions~\cite{who2024}. Recent advances in computational pathology
and pathology foundation models~\cite{chen2024uni, vorontsov2024virchow,
lu2024conch} have enabled integration of complementary
modalities---particularly whole-slide histopathology images and bulk
RNA-seq transcriptomes~\cite{chen2021, chen2022porpoise,
acosta2022}---yielding prognostic signal that neither modality delivers
alone.

Joint training of image and genomic encoders under Cox partial
likelihood~\cite{cox1972} faces two compounding challenges.
First, \textit{modality imbalance}: one encoder's gradients often
dominate the shared optimization, leaving the other modality
under-trained~\cite{wang2020what, peng2022ogm}. Second,
\textit{objective conflict}: auxiliary training objectives---such as
per-modality survival heads or pairwise ranking
losses---can produce gradients that \textit{oppose} the primary
objective's descent direction in encoder parameter space. While
multi-task learning and gradient modulation
methods~\cite{peng2022ogm, du2023unimodal, fan2023pmr} address
modality imbalance, they do not explicitly resolve inter-objective
gradient conflicts.

We introduce \textbf{RA-SAMO} (Representation-Aware Survival-Aware
Multi-Objective optimization), a framework that addresses both
challenges jointly. RA-SAMO trains a multimodal survival model with a
primary Cox loss and multiple auxiliary objectives (per-modality Cox
heads, pairwise ranking loss), adaptively weights them using an
EMA-based controller, and---critically---projects auxiliary gradients
that conflict with the primary objective onto its orthogonal
complement before the optimizer step. This ensures the primary
objective's descent direction is never opposed while retaining
non-conflicting auxiliary signal.

Through extensive evaluation on five TCGA cancer cohorts and the
HAM10000 dermoscopic classification benchmark, we make three
contributions:

\begin{itemize}
\item \textbf{Gradient projection for multi-objective medical
learning.} A conflict-resolution mechanism that projects conflicting
auxiliary encoder gradients onto the orthogonal complement of the
primary gradient, preventing auxiliary objectives from degrading
the primary loss (Section~\ref{sec:projection}).

\item \textbf{Comprehensive gradient-level diagnostics.} Fisher
information traces and inter-objective cosine similarities reveal
that objective utility is cancer-dependent: ranking loss is
near-redundant with Cox ($\cos > 0.93$), while auxiliary survival
heads provide orthogonal but noisy signal ($\cos \approx 0.14$)
(Section~\ref{sec:diagnostics}).

\item \textbf{Cross-domain generalization.} RA-SAMO generalizes
from survival prediction to 7-class classification on HAM10000
(image + tabular metadata), producing the most stable performance
across random seeds (Section~\ref{sec:ham10000}).
\end{itemize}


% ─────────────────────────────────────────────────────────────────────────────
\section{Related Work}

\subsection{Multimodal Cancer Survival Prediction}

MCAT~\cite{chen2021} uses co-attention to fuse gigapixel slides with
genomic features. PORPOISE~\cite{chen2022porpoise} performs pan-cancer
integrative analysis via multimodal deep learning.
SurvPath~\cite{jaume2024survpath} models dense interactions between
biological pathways and histology. MOTCat~\cite{xu2023motcat}
introduces optimal transport-based co-attention for cross-modal
alignment. These architectures assume standard optimizers suffice and
do not address gradient conflicts between training objectives.

\subsection{Modality Imbalance}

Wang et al.~\cite{wang2020what} identified that one modality often
dominates multimodal training. OGM-GE~\cite{peng2022ogm} suppresses
dominant-modality gradients with an adaptive scaling factor and
Gaussian noise injection. PMR~\cite{fan2023pmr} introduces prototypical
modal rebalancing. Du et al.~\cite{du2023unimodal} use unimodal
teachers for balanced learning. These methods operate on gradient
\textit{magnitudes} but do not consider gradient \textit{directions}
between objectives.

\subsection{Multi-Objective Gradient Methods}

Multi-task learning faces conflicting gradients across tasks. PCGrad
\cite{yu2020pcgrad} projects conflicting task gradients; CAGrad
\cite{liu2021cagrad} finds a common descent direction; Nash-MTL
\cite{navon2022nashmtl} frames it as a bargaining game. These methods
resolve inter-task conflicts but have not been applied to
multi-objective training within a single multimodal survival model,
where the ``tasks'' are different loss functions on the same
prediction target.

\subsection{Sharpness-Aware Minimization}

SAM~\cite{foret2021sam} seeks flat minima via worst-case perturbation.
M-SAM~\cite{rajoli2025msam} targets perturbation to the dominant
modality. RA-SAMO is complementary: it can wrap any base optimizer
(including SAM) and adds conflict resolution at the objective level
rather than the perturbation level.


% ─────────────────────────────────────────────────────────────────────────────
\section{Method}

\subsection{Problem Formulation}

Given $N$ patients, each with histopathology slide embedding
$\mathbf{x}_i^{\text{img}} \in \mathbb{R}^{d_{\text{img}}}$ and gene
expression profile
$\mathbf{x}_i^{\text{gene}} \in \mathbb{R}^{d_{\text{gene}}}$, with
survival time $t_i > 0$ and event indicator $\delta_i \in \{0, 1\}$
($1 =$ observed death), we seek a risk function $f_\theta$ such that
patients with higher predicted risk
$r_i = f_\theta(\mathbf{x}_i^{\text{img}},
\mathbf{x}_i^{\text{gene}})$ have shorter survival times.
Performance is measured by Harrell's concordance index
(C-index)~\cite{harrell1996}.

\subsection{Model Architecture}

\textbf{Image encoder.}
Pre-extracted slide embeddings from UNI~\cite{chen2024uni} (1024-d for
BRCA) or UNI2-h (1536-d for others) are mapped via a single-layer MLP:
\begin{equation}
\mathbf{h}^{\text{img}} = \text{Dropout}_{p}\!\bigl(
\text{ReLU}\bigl(\text{BN}(\mathbf{W}_{\text{img}} \mathbf{x}^{\text{img}})\bigr)\bigr)
\in \mathbb{R}^{512}
\label{eq:img_enc}
\end{equation}

\textbf{Gene encoder.}
A three-layer MLP encodes gene expression:
\begin{equation}
\mathbf{h}^{\text{gene}} = \text{MLP}_{512 \to 256 \to 128}(\mathbf{x}^{\text{gene}})
\in \mathbb{R}^{128}
\label{eq:gene_enc}
\end{equation}
with batch normalization, ReLU, and dropout ($p\!=\!0.3$) after each layer.

\textbf{Cross-modal attention fusion.}
Both representations are projected to a shared $d\!=\!256$ space, and a
learned attention network computes per-patient modality weights:
\begin{equation}
\alpha_m = \frac{
  \exp\!\bigl(\mathbf{u}^\top \tanh(\mathbf{W}_a \tilde{\mathbf{h}}^m)\bigr)
}{
  \sum_{m'} \exp\!\bigl(\mathbf{u}^\top \tanh(\mathbf{W}_a \tilde{\mathbf{h}}^{m'})\bigr)
}
\label{eq:attn}
\end{equation}
producing the fused representation
$\mathbf{h} = \sum_m \alpha_m \tilde{\mathbf{h}}^m \in \mathbb{R}^{256}$.

\textbf{Survival head.}
A shared layer ($256 \to 256$, ReLU, dropout) followed by a linear
projection yields risk score $r_i$.

\textbf{Auxiliary heads.}
Per-modality auxiliary Cox survival heads
($\mathbb{R}^{512} \to \mathbb{R}^{64} \to \mathbb{R}^{1}$ for image;
$\mathbb{R}^{128} \to \mathbb{R}^{64} \to \mathbb{R}^{1}$ for gene)
provide unimodal survival predictions used as auxiliary objectives.

\subsection{Multi-Objective Loss}

The total loss combines four components:
\begin{equation}
\mathcal{L} = \mathcal{L}_{\text{cox}}^{\text{fuse}} +
w_{\text{img}} \mathcal{L}_{\text{cox}}^{\text{img}} +
w_{\text{gene}} \mathcal{L}_{\text{cox}}^{\text{gene}} +
w_{\text{rank}} \mathcal{L}_{\text{rank}}
\label{eq:total_loss}
\end{equation}

\textbf{Primary loss.} The Cox partial log-likelihood~\cite{cox1972}
on fused predictions with Breslow's tied-event approximation:
\begin{equation}
\mathcal{L}_{\text{cox}} = -\frac{1}{|\mathcal{E}|}
\sum_{i: \delta_i = 1}
\left( r_i - \log \sum_{j: t_j \geq t_i} e^{r_j} \right)
\label{eq:cox}
\end{equation}
where $\mathcal{E} = \{i : \delta_i = 1\}$ is the set of observed events.

\textbf{Auxiliary Cox losses.}
$\mathcal{L}_{\text{cox}}^{\text{img}}$ and
$\mathcal{L}_{\text{cox}}^{\text{gene}}$ apply Eq.~\ref{eq:cox} to
per-modality risk scores from auxiliary heads. These encourage each
encoder to independently learn prognostically informative
representations.

\textbf{Pairwise ranking loss.}
For comparable pairs where patient $i$ has an observed event before
patient $j$'s follow-up ends:
\begin{equation}
\mathcal{L}_{\text{rank}} = \frac{1}{|\mathcal{P}|}
\sum_{(i,j) \in \mathcal{P}}
\log\bigl(1 + e^{(r_j - r_i)/\tau}\bigr)
\label{eq:rank}
\end{equation}
where $\mathcal{P} = \{(i,j) : \delta_i = 1 \land t_i < t_j\}$ and
$\tau$ is a temperature parameter.

\subsection{Adaptive Loss Weighting}

The weights $\mathbf{w} = (w_{\text{img}}, w_{\text{gene}},
w_{\text{rank}})$ are updated per epoch using an EMA-based controller.
Let $\bar{\ell}_k^{(e)}$ be the exponentially smoothed epoch-mean loss
for component $k$ at epoch $e$, with smoothing factor $\beta = 0.9$:
\begin{equation}
\bar{\ell}_k^{(e)} = \beta \cdot \bar{\ell}_k^{(e-1)} + (1 - \beta) \cdot \ell_k^{(e)}
\end{equation}

The controller computes normalized ratios against first-epoch baselines
and adjusts weights toward components whose losses are lagging:
\begin{equation}
w_k^{(e+1)} = \text{clip}\!\left(
w_k^{(e)} + \text{clip}\!\left(
\frac{0.1 \cdot \bar{\ell}_k^{(e)} / \bar{\ell}_k^{(0)}}
{\text{mean}_j(\bar{\ell}_j^{(e)} / \bar{\ell}_j^{(0)})} - w_k^{(e)},
\pm \Delta_{\max}
\right),
[w_{\min}, w_{\max}]
\right)
\label{eq:weight_update}
\end{equation}
with $w_{\min}\!=\!0.02$, $w_{\max}\!=\!0.3$,
$\Delta_{\max}\!=\!0.02$, and a warmup of 2 epochs during which
weights remain at their initial values (0.1 each). This bounded,
slow-moving controller prevents any auxiliary objective from
overwhelming the primary loss.

\subsection{Gradient Projection for Conflict Resolution}
\label{sec:projection}

The key contribution of RA-SAMO is explicit resolution of gradient
conflicts between the primary and auxiliary objectives in each
encoder's parameter space.

For each encoder $m \in \{\text{img}, \text{gene}\}$, let
$\mathbf{g}_m^{\text{pri}}$, $\mathbf{g}_m^{\text{aux}}$, and
$\mathbf{g}_m^{\text{rank}}$ denote the gradients of the primary
(fused Cox), auxiliary (per-modality Cox), and ranking losses with
respect to encoder $m$'s parameters $\theta_m$, computed via
\texttt{torch.autograd.grad} with \texttt{retain\_graph=True}.

\textbf{Conflict detection.} An auxiliary gradient
$\mathbf{g}_m^{\text{sec}}$ (where ``sec'' is either aux or rank)
\textit{conflicts} with the primary if their inner product is negative:
\begin{equation}
\text{conflict}(\mathbf{g}_m^{\text{sec}}, \mathbf{g}_m^{\text{pri}}) =
\bigl(\langle \mathbf{g}_m^{\text{sec}}, \mathbf{g}_m^{\text{pri}} \rangle < 0\bigr)
\label{eq:conflict}
\end{equation}

\textbf{Projection.} When conflict is detected, the auxiliary gradient
is projected onto the orthogonal complement of the primary gradient:
\begin{equation}
\hat{\mathbf{g}}_m^{\text{sec}} = \mathbf{g}_m^{\text{sec}} -
\frac{\langle \mathbf{g}_m^{\text{sec}}, \mathbf{g}_m^{\text{pri}} \rangle}
{\|\mathbf{g}_m^{\text{pri}}\|^2}
\mathbf{g}_m^{\text{pri}}
\label{eq:project}
\end{equation}
When no conflict exists ($\langle \mathbf{g}_m^{\text{sec}},
\mathbf{g}_m^{\text{pri}} \rangle \geq 0$), the auxiliary gradient is
used unmodified.

\textbf{Gradient replacement.} The final encoder gradient combines
primary and (possibly projected) auxiliary gradients:
\begin{equation}
\mathbf{g}_m = \mathbf{g}_m^{\text{pri}} +
w_{\text{aux}} \hat{\mathbf{g}}_m^{\text{aux}} +
w_{\text{rank}} \hat{\mathbf{g}}_m^{\text{rank}}
\label{eq:combined}
\end{equation}
This replaces the encoder's \texttt{.grad} attribute after
\texttt{backward()} and before \texttt{optimizer.step()}.
Shared parameters (fusion, survival head) receive standard
gradients from the weighted loss.

The full training step is summarized in Algorithm~\ref{alg:rasamo}.

\begin{algorithm}[t]
\caption{RA-SAMO Training Step}
\label{alg:rasamo}
\begin{algorithmic}[1]
\REQUIRE Model $f_\theta$ with encoders $\theta_{\text{img}}, \theta_{\text{gene}}$,
batch $(\mathbf{X}, \mathbf{t}, \boldsymbol{\delta})$,
weights $(w_{\text{img}}, w_{\text{gene}}, w_{\text{rank}})$
\STATE Forward pass: compute
$\mathcal{L}_{\text{cox}}^{\text{fuse}}$,
$\mathcal{L}_{\text{cox}}^{\text{img}}$,
$\mathcal{L}_{\text{cox}}^{\text{gene}}$,
$\mathcal{L}_{\text{rank}}$
\FOR{each encoder $m \in \{\text{img}, \text{gene}\}$}
  \STATE $\mathbf{g}_m^{\text{pri}} \leftarrow \nabla_{\theta_m} \mathcal{L}_{\text{cox}}^{\text{fuse}}$
  \STATE $\mathbf{g}_m^{\text{aux}} \leftarrow \nabla_{\theta_m} \mathcal{L}_{\text{cox}}^{m}$
  \STATE $\mathbf{g}_m^{\text{rank}} \leftarrow \nabla_{\theta_m} \mathcal{L}_{\text{rank}}$
  \IF{$\langle \mathbf{g}_m^{\text{aux}}, \mathbf{g}_m^{\text{pri}} \rangle < 0$}
    \STATE Project: $\hat{\mathbf{g}}_m^{\text{aux}} \leftarrow$ Eq.~\ref{eq:project}
  \ELSE
    \STATE $\hat{\mathbf{g}}_m^{\text{aux}} \leftarrow \mathbf{g}_m^{\text{aux}}$
  \ENDIF
  \STATE Repeat projection for $\mathbf{g}_m^{\text{rank}}$
  \STATE $\mathbf{g}_m \leftarrow \mathbf{g}_m^{\text{pri}} + w_{\text{aux}} \hat{\mathbf{g}}_m^{\text{aux}} + w_{\text{rank}} \hat{\mathbf{g}}_m^{\text{rank}}$
\ENDFOR
\STATE $\mathcal{L}_{\text{total}} \leftarrow$ Eq.~\ref{eq:total_loss}
\STATE $\mathcal{L}_{\text{total}}$.\texttt{backward()}
\STATE Replace encoder gradients: $\theta_m$.\texttt{grad} $\leftarrow \mathbf{g}_m$
\STATE \texttt{optimizer.step()}
\STATE Update weights $\mathbf{w}$ via Eq.~\ref{eq:weight_update} (end of epoch)
\end{algorithmic}
\end{algorithm}

\textbf{Diagnostic output.} At each step, RA-SAMO records
per-encoder cosine similarities
$\cos(\mathbf{g}_m^{\text{aux}}, \mathbf{g}_m^{\text{pri}})$ and
$\cos(\mathbf{g}_m^{\text{rank}}, \mathbf{g}_m^{\text{pri}})$,
whether projection was applied, and the relative update change
$\|\mathbf{g}_m^{\text{proj}} - \mathbf{g}_m^{\text{orig}}\| /
\|\mathbf{g}_m^{\text{orig}}\|$. These diagnostics enable
post-hoc analysis of when and where gradient conflicts occur
(Section~\ref{sec:diagnostics}).


% ─────────────────────────────────────────────────────────────────────────────
\section{Experimental Setup}

\subsection{Datasets}

\textbf{TCGA Survival.}
We evaluate on five TCGA cancer cohorts (Table~\ref{tab:datasets}).
Histopathology slide embeddings are extracted using
UNI~\cite{chen2024uni} (1024-d for BRCA) or UNI2-h (1536-d for
others), following the pre-extracted embeddings released by
Jaume et al.~\cite{jaume2024tangle}. Gene expression profiles use
the top 1,000 most variable genes per cohort after log1p
transformation and standard scaling (fitted on training data only).

\begin{table}[t]
\centering
\caption{Dataset statistics. Event rate = fraction of patients with
observed death (uncensored).}
\label{tab:datasets}
\begin{tabular}{lrrrl}
\toprule
Cancer & $N$ & Event \% & Cens. \% & Encoder \\
\midrule
BRCA & 1,123 & 13.6\% & 86.4\% & UNI 1024-d \\
KIRC & 519   & 34.6\% & 65.4\% & UNI2-h 1536-d \\
LUAD & 522   & 35.6\% & 64.4\% & UNI2-h 1536-d \\
LUSC & 500   & 42.6\% & 57.4\% & UNI2-h 1536-d \\
STAD & 367   & 39.4\% & 60.6\% & UNI2-h 1536-d \\
\bottomrule
\end{tabular}
\end{table}

\textbf{HAM10000 Classification.}
To test generalization beyond survival prediction, we evaluate on
HAM10000~\cite{tschandl2018ham10000}: 10,015 dermoscopic images with
tabular metadata (age, sex, localization) across 7 lesion classes.
Splits are grouped by lesion ID to prevent data leakage (same
lesion may have multiple images), with stratified 70/15/15
train/val/test splits.

\subsection{Methods Compared}

All methods use the same architecture and hyperparameters; only the
optimization strategy varies:

\begin{itemize}
\item \textbf{Baseline}: Cox loss only (Eq.~\ref{eq:cox}), AdamW optimizer.
\item \textbf{+ Aux}: Add per-modality auxiliary Cox heads.
\item \textbf{+ Rank}: Add pairwise ranking loss (Eq.~\ref{eq:rank}).
\item \textbf{+ Aux+Rank}: Both auxiliary heads and ranking loss.
\item \textbf{Adaptive Aux}: Auxiliary heads with EMA-based weight
controller (Eq.~\ref{eq:weight_update}), no gradient projection.
\item \textbf{SAMO}: Adaptive Aux + Rank with weight controller.
\item \textbf{RA-SAMO}: SAMO + gradient projection
(Eq.~\ref{eq:project}--\ref{eq:combined}).
\end{itemize}

For HAM10000, we additionally compare SAM~\cite{foret2021sam} and
unimodal baselines (image-only, metadata-only), replacing Cox loss
with cross-entropy.

\subsection{Training Protocol}

\textbf{TCGA:} 50 epochs, batch size 32, AdamW (lr$\!=\!10^{-3}$,
weight decay$\!=\!10^{-4}$), cosine annealing. Stratified 5-fold
cross-validation (stratified by vital status). Model selection on
best validation C-index. Each patient tested exactly once.

\textbf{HAM10000:} 30 epochs, batch size 32, AdamW
(lr$\!=\!10^{-3}$), cosine annealing. ResNet-18 (pretrained
ImageNet) image encoder, 3-layer MLP metadata encoder. Three seeds
(42, 43, 44). Model selection on best validation macro-F1. Auxiliary
objectives use cross-entropy instead of Cox. Ranking loss weight
set to zero (no natural ordering in classification).

All experiments run on a single NVIDIA T4 GPU (16 GB).


% ─────────────────────────────────────────────────────────────────────────────
\section{Results}

\subsection{TCGA Survival Prediction}

Table~\ref{tab:tcga_results} presents test C-index across five TCGA
cohorts under 5-fold cross-validation.

\begin{table*}[t]
\centering
\caption{Test C-index (5-fold mean $\pm$ std) on TCGA survival
prediction. Best per-cancer in \textbf{bold}. ``Adap.~Aux'' =
Adaptive Aux. All methods use the same architecture and hyperparameters.}
\label{tab:tcga_results}
\begin{tabular}{lcccccr}
\toprule
Method & BRCA & KIRC & LUAD & LUSC & STAD & Mean \\
\midrule
Baseline (Cox) &
  $.629 \pm .067$ &
  $\mathbf{.740} \pm .061$ &
  $.641 \pm .058$ &
  $.559 \pm .052$ &
  $.583 \pm .023$ &
  .630 \\
+ Aux &
  $.662 \pm .056$ &
  $.728 \pm .062$ &
  $.618 \pm .041$ &
  $.540 \pm .066$ &
  $.581 \pm .033$ &
  .626 \\
+ Rank &
  $.621 \pm .056$ &
  $.731 \pm .062$ &
  $\mathbf{.646} \pm .041$ &
  $\mathbf{.568} \pm .066$ &
  $.591 \pm .033$ &
  .631 \\
+ Aux+Rank &
  $.650 \pm .056$ &
  $.731 \pm .062$ &
  $.637 \pm .041$ &
  $.551 \pm .066$ &
  $.568 \pm .033$ &
  .628 \\
Adap.\ Aux &
  $\mathbf{.678} \pm .066$ &
  $.731 \pm .062$ &
  $.626 \pm .041$ &
  $.539 \pm .071$ &
  $\mathbf{.595} \pm .038$ &
  $\mathbf{.634}$ \\
SAMO &
  $.632 \pm .056$ &
  $.731 \pm .062$ &
  $.631 \pm .041$ &
  $.536 \pm .066$ &
  $.569 \pm .033$ &
  .620 \\
RA-SAMO &
  $.640 \pm .067$ &
  $.729 \pm .061$ &
  $.629 \pm .058$ &
  $.562 \pm .052$ &
  $.583 \pm .023$ &
  .629 \\
\bottomrule
\end{tabular}
\end{table*}

\textbf{No single method dominates.} The best method is
cancer-dependent: Adaptive Aux wins on BRCA and STAD, pure ranking
wins on LUAD and LUSC, and the Cox-only baseline wins on KIRC. All
differences between the top-ranked methods are within noise
($< 0.005$ mean C-index).

\textbf{RA-SAMO is consistently competitive.} While not the top
performer on any single cancer, RA-SAMO (mean 0.629) is within
0.005 of the best method overall (Adaptive Aux, 0.634) and notably
improves over SAMO (0.620) by 0.009. The gradient projection
prevents the degradation observed in SAMO, where conflicting
auxiliary gradients in the ranking branch hurt performance.

\textbf{Ranking loss is redundant with Cox.} The ``+ Rank'' method
achieves high C-index on LUAD/LUSC but not BRCA, consistent with
gradient alignment analysis (Section~\ref{sec:diagnostics}): ranking
gradients are $>93\%$ aligned with Cox, adding signal only when the
ranking pairs provide complementary information.

\subsection{HAM10000 Classification}
\label{sec:ham10000}

Table~\ref{tab:ham10000} shows results on 7-class dermoscopic
classification, demonstrating that RA-SAMO's multi-objective
framework generalizes beyond survival prediction.

\begin{table}[t]
\centering
\caption{HAM10000 7-class classification (mean $\pm$ std over 3
seeds). Best in \textbf{bold}; lowest std in \underline{underline}.}
\label{tab:ham10000}
\begin{tabular}{lcc}
\toprule
Method & Macro-F1 & Balanced Acc \\
\midrule
Image-Only & $.698 \pm .018$ & $.670 \pm .020$ \\
Meta-Only  & $.196 \pm .008$ & $.216 \pm .016$ \\
AdamW      & $\mathbf{.709} \pm .028$ & $\mathbf{.704} \pm .028$ \\
SAM        & $.704 \pm .033$ & $.683 \pm .035$ \\
Adap.\ Aux & $.692 \pm .028$ & $.670 \pm .034$ \\
SAMO       & $.685 \pm .027$ & $.661 \pm .040$ \\
RA-SAMO    & $\underline{.695 \pm .011}$ & $.667 \pm .023$ \\
\bottomrule
\end{tabular}
\end{table}

\textbf{RA-SAMO has the lowest variance.} While AdamW achieves the
highest absolute macro-F1 (0.709), its cross-seed standard
deviation is 2.5$\times$ higher than RA-SAMO's (0.028 vs 0.011).
RA-SAMO's gradient projection smooths the optimization landscape,
producing more consistent results across random initializations.

\textbf{Metadata contributes minimally.} The Meta-Only baseline
achieves only 0.196 macro-F1, confirming that the image modality
dominates. This creates a natural testbed for multimodal optimization:
methods must learn to utilize the weaker modality without being
overwhelmed by the stronger one.


% ─────────────────────────────────────────────────────────────────────────────
\section{Gradient-Level Diagnostics}
\label{sec:diagnostics}

We conducted a 27-run diagnostic study (3 cancers $\times$ 3
methods $\times$ 3 folds $\times$ 50 epochs) recording Fisher
information traces and inter-objective gradient alignments at
diagnostic epochs.

\subsection{Inter-Objective Gradient Alignment}

Table~\ref{tab:alignment} reports cosine similarity between the
primary Cox gradient and auxiliary objective gradients, measured in
encoder parameter space.

\begin{table}[t]
\centering
\caption{Gradient cosine similarity: primary Cox vs.\ secondary
objectives (mean $\pm$ std across folds and diagnostic epochs).}
\label{tab:alignment}
\begin{tabular}{lcc}
\toprule
Cancer & $\cos(\nabla\mathcal{L}_{\text{cox}}, \nabla\mathcal{L}_{\text{aux}})$ &
$\cos(\nabla\mathcal{L}_{\text{cox}}, \nabla\mathcal{L}_{\text{rank}})$ \\
\midrule
BRCA & $0.180 \pm 0.058$ & $0.981 \pm 0.018$ \\
KIRC & $0.168 \pm 0.074$ & $0.947 \pm 0.025$ \\
LUAD & $0.109 \pm 0.062$ & $0.947 \pm 0.027$ \\
LUSC & $0.106 \pm 0.053$ & $0.933 \pm 0.031$ \\
STAD & $0.121 \pm 0.055$ & $0.938 \pm 0.029$ \\
\bottomrule
\end{tabular}
\end{table}

\textbf{Ranking loss is near-redundant.} Across all five cancers,
$\cos(\nabla\mathcal{L}_{\text{cox}}, \nabla\mathcal{L}_{\text{rank}})
> 0.93$. The ranking loss pushes parameters in essentially the same
direction as Cox, providing minimal orthogonal information. This
explains why ``+ Rank'' alone rarely improves over the baseline.

\textbf{Auxiliary heads provide orthogonal signal.} The cosine
similarity between Cox and auxiliary gradients is low
($0.11$--$0.18$), indicating that per-modality survival objectives
provide genuinely different gradient directions. However, this
orthogonality means auxiliary gradients may also \textit{conflict}
with the primary objective, motivating RA-SAMO's projection mechanism.

\subsection{Fisher Information and Modality Imbalance}

The Fisher information trace
$\text{Tr}(\hat{F}_m) = \frac{1}{|\mathcal{E}|} \sum_{i \in \mathcal{E}}
\|\nabla_{\theta_m} \ell_i^{\text{Cox}}\|^2$
quantifies how much gradient signal each encoder receives from
observed events.

\textbf{Gene encoder dominates.} The gene-to-image Fisher ratio
exceeds 1.0 across all cancers and methods: 1.12--1.54 for BRCA,
1.99--2.40 for KIRC, 2.21--2.39 for LUAD. This indicates the gene
encoder consistently receives stronger per-event gradient signal,
contrary to the image dominance reported in audiovisual
settings~\cite{peng2022ogm}. In the survival prediction setting with
pre-extracted image features, the gene encoder has more trainable
parameters and larger gradient magnitudes.

\textbf{Fisher imbalance does not predict intervention benefit.}
We tested whether early-epoch (epochs 1--5) Fisher ratio predicts
test C-index improvement from auxiliary objectives. The Pearson
correlation is $r = -0.36$, $p = 0.15$---not statistically
significant. This suggests that modality imbalance alone does not
determine whether multi-objective training will help; the gradient
\textit{direction} (conflict vs.\ alignment) matters more than
magnitude.


% ─────────────────────────────────────────────────────────────────────────────
\section{Discussion}

\subsection{Cancer-Dependent Objective Selection}

The most striking finding is that no single objective combination
dominates across all cancers. Table~\ref{tab:tcga_results} shows:
\begin{itemize}
\item \textbf{BRCA} (86\% censored): Adaptive Aux wins (0.678).
Heavy censoring makes auxiliary per-modality heads valuable for
representation learning.
\item \textbf{KIRC} (65\% censored): Baseline Cox wins (0.740).
Moderate censoring and relatively balanced modalities mean auxiliary
objectives add noise.
\item \textbf{LUAD/LUSC}: Ranking loss wins. Moderate event rates
provide enough comparable pairs for the ranking objective to add signal.
\item \textbf{STAD} (61\% censored): Adaptive Aux wins (0.595).
Small cohort size ($n\!=\!367$) benefits from the regularizing
effect of auxiliary heads.
\end{itemize}

This cancer-dependence has practical implications: rather than
applying a single optimizer to all cohorts, practitioners should
validate objective combinations per cancer type.

\subsection{When Does Gradient Projection Help?}

RA-SAMO improves over SAMO (which lacks projection) by 0.009 mean
C-index. The projection is most valuable when:
\begin{enumerate}
\item Auxiliary gradients are both \textit{informative} (providing
signal not captured by the primary loss) and \textit{conflicting}
(opposing the primary gradient direction).
\item The auxiliary weight is large enough for conflicts to materially
affect the optimizer step.
\end{enumerate}

When auxiliary gradients are aligned with the primary (as ranking loss
nearly always is), projection is rarely triggered and has no effect.
When auxiliary gradients are orthogonal but non-conflicting, they
pass through unmodified. The projection mechanism specifically targets
the harmful case: strong auxiliary signal pushing \textit{against}
the primary objective.

\subsection{Limitations and Future Work}

\textbf{Scale of differences.} All methods in Table~\ref{tab:tcga_results}
fall within a narrow C-index band (0.620--0.634 mean). With 5-fold
cross-validation on cohorts of 367--1,123 patients, individual fold
variance is large (std 0.02--0.07). Larger cohorts or more folds
are needed for definitive statistical separation.

\textbf{Fixed architecture.} Results are specific to the late-fusion
attention architecture with pre-extracted features. Early-fusion
architectures or end-to-end training may show different gradient
dynamics.

\textbf{Loss-progress vs.\ performance.} Our adaptive weight
controller uses loss-progress signals (Eq.~\ref{eq:weight_update}),
which are not always correlated with C-index improvement. A
controller that directly optimizes validation C-index (e.g., via
short inner-validation rollouts) may provide better objective
selection, though our experiments with 2-epoch rollouts
showed insufficient differential signal.

\textbf{Computational cost.} RA-SAMO requires per-encoder gradient
computation via \texttt{autograd.grad} (6 calls per step: 2
encoders $\times$ 3 objectives), adding $\sim$10--15\% wall-clock
overhead versus standard backpropagation. This is modest compared
to SAM's 2$\times$ forward-backward cost.


% ─────────────────────────────────────────────────────────────────────────────
\section{Conclusion}

We introduced RA-SAMO, a multi-objective optimization framework for
multimodal medical learning that resolves inter-objective gradient
conflicts via orthogonal projection. Evaluation across five TCGA
cancer cohorts and the HAM10000 classification benchmark reveals
three key findings: (1)~optimal objective combinations are
cancer-dependent, with no single method dominating; (2)~ranking
loss is near-redundant with Cox partial likelihood ($\cos > 0.93$),
while per-modality auxiliary heads provide genuinely orthogonal
signal; (3)~gradient projection prevents the degradation observed
in na\"ive multi-objective training, yielding the most consistent
performance across cancers and random seeds.

These results suggest that optimizer design---specifically, how
multiple objectives interact in the gradient space of each
encoder---deserves systematic attention alongside architecture
and representation choices in multimodal medical AI.


% ─────────────────────────────────────────────────────────────────────────────
\bibliographystyle{plainnat}
\bibliography{references}

\end{document}
